% Generated draft: requires mathematical and source review.
\documentclass[11pt]{article}
\usepackage[margin=25mm]{geometry}
\usepackage{fontspec}
\setmainfont{Latin Modern Roman}
\usepackage{amsmath,amssymb,mathtools,unicode-math}
\setmathfont{Latin Modern Math}
\usepackage{enumitem,needspace,xurl,hyperref}
\setlength{\parindent}{0pt}
\setlength{\parskip}{5pt}
\setlength{\emergencystretch}{3em}
\newcommand{\N}{\mathbb{N}}
\newcommand{\Z}{\mathbb{Z}}
\newcommand{\Q}{\mathbb{Q}}
\newcommand{\R}{\mathbb{R}}
\newcommand{\C}{\mathbb{C}}
\newcommand{\F}{\mathbb{F}}
\newcommand{\A}{\mathbb{A}}
\newcommand{\PP}{\mathbb{P}}
\newcommand{\E}{\mathbb{E}}
\newcommand{\eps}{\varepsilon}
\newcommand{\dd}{\,\mathrm{d}}
\begin{document}
\section*{UnsolvedMath ID: 10900047}
\textbf{Source problem code:} catalogue entry
\subsection{Definitions and mathematical statement}
Maher's Problem 6.3 asks to investigate the apparent extreme regularity of Rivin's experimental results and the possibility that it reflects additional mathematical structure. The source supplies no particular experiment, dataset, observable, sampling distribution, parameter, numerical regularity or proposed law.

The supplied research summary mentions experiments on random triangulations of \(3\)-manifolds and counts of embedded essential surfaces, and related work on mapping-class-group random walks and random triangulations. These descriptions do not identify a single probability model or theorem. A precise investigation would first have to fix the experimental ensemble, measured quantities and meaning of regularity before explaining a pattern or formulating a statistical or structural claim. No limiting distribution, convergence rate, exact identity or hidden structure is asserted by the record. Its word ``possibly'' makes the structural interpretation a hypothesis to investigate, not an established assumption.

The summary reports that work by Maher, Dunfield and others explains some observed regularity but that no verified complete structural explanation of all of Rivin's experiments is available. This broad assessment has no supplied references or defined list of experiments, and is retained without elevating it into a formal nonexistence or exhaustive open-status theorem.
\subsection{Short English statement}
Investigate why the referenced Rivin experiments look unusually regular and whether an underlying structure explains them. The experiments and the regularity are not specified, so a mathematical target must be recovered from context. The supplied summary reports partial explanations but no verified complete one.
\subsection{Sources}
\begin{itemize}
\item Local UnsolvedMath registry, permanent ID \(10900047\), exact source code \texttt{catalogue entry}, subsection 6.3, attributed to Maher; no proposed year supplied. The experiments are attributed to Rivin.
\item Delp, Hoffoss and Manning, \emph{Problems In Groups, Geometry, and Three-Manifolds} (2015), 6.3, PDF page 5: \url{https://arxiv.org/pdf/1512.04620#page=5}; supplied status-source URL \url{https://arxiv.org/abs/1512.04620}. No source-side resolution is recorded.
\item Supplied summary names Maher, Dunfield and others for related random-walk/triangulation work but supplies no complete references or experimental data. No separate authoritative citation identifying the experiments is provided.
\item Source-TeX extraction checked against rendered PDF text, accessed 2026-07-29. Source evidence NEEDS\_REVIEW: open in 2015, with no complete modern primary-literature audit recorded. Registry status open and classification OPEN-TRIAGE; summary calls the complete explanatory question unresolved.
\item Dataset attribution: Ulam AI / UnsolvedMath Contributors, CC BY 4.0. Source rights NEEDS\_REVIEW: publicly accessible author-submitted source TeX and PDF, redistribution terms requiring release review.
\end{itemize}
\end{document}
